\section{Timing}

\subsection{Combinational Timing}

\begin{definition}
\textbf{Propagation Delay ($t_{pd}$)}: Maximum time from input change until output reaches final value.
\textbf{Contamination Delay ($t_{cd}$)}: Minimum time from input change until output starts to change.
\end{definition}

\begin{example}
\textbf{Glitch}: Single input transition causing multiple output transitions due to different path delays. Fix with redundant "consensus terms" in K-maps.
\end{example}

\subsection{Sequential Timing}

\begin{definition}
\textbf{Aperture time} ($t_a = t_{setup} + t_{hold}$):
\begin{itemize}
    \item \textbf{Setup Time ($t_{setup}$)}: Time before clock edge data must be stable.
    \item \textbf{Hold Time ($t_{hold}$)}: Time after clock edge data must remain stable.
\end{itemize}
Violations lead to \textbf{metastable} states.
\end{definition}

\begin{definition}
\textbf{Prop. Clock-to-Q Delay ($t_{pcq}$)}: Max time after clock edge until $Q$ is stable.
\textbf{Cont. Clock-to-Q Delay ($t_{ccq}$)}: Min time after clock edge until $Q$ starts to change.
\end{definition}

\begin{theorem}
For a synchronous system to operate correctly:
\textbf{Setup Constraint (Max Freq)}: $$T_c \geq t_{pcq} + t_{pd} + t_{setup}$$
\textbf{Hold Constraint}: $$t_{ccq} + t_{cd} > t_{hold}$$
(Hold time violations cannot be fixed by simply slowing down the clock).
\end{theorem}

\subsection{Clock Skew}

\begin{definition}
\textbf{Clock skew} ($t_{skew}$): Time difference between arrival of clock signal at different flip-flops. Worsens both constraints:
\begin{itemize}
    \item \textbf{Setup}: $T_c \geq t_{pcq} + t_{pd} + t_{setup} + t_{skew}$
    \item \textbf{Hold}: $t_{ccq} + t_{cd} > t_{hold} + t_{skew}$
\end{itemize}
\end{definition}
